% =====================================================================
% 03_architecture.tex
% =====================================================================
\begingroup
\newcolumntype{L}[1]{>{\hsize=#1\hsize\raggedright\arraybackslash}X}

\section{Kiến trúc và các thành phần WLAN}
\label{sec:arch}

\subsection{Các thành phần cơ bản}

\begin{description}[style=nextline]
  \item[\acl{sta}] Thiết bị có bộ thu phát vô tuyến 802.11, chẳng hạn máy tính
    xách tay, điện thoại thông minh, máy tính để bàn có card mạng không dây, cảm biến \ac{iot} hoặc máy in \cite{kurose2012, cisco_intro_80211, microsoft_wifi}. Trạm khách \ac{sta}
    thiết lập liên kết với một \ac{ap}.
  \item[\acl{ap}] Cầu nối các thiết bị khách không dây với \ac{lan} có dây \cite{kurose2012, cisco_intro_80211}. Một
    \ac{ap} có thể phát nhiều mạng ảo, mỗi mạng có một \ac{ssid} riêng.
  \item[\acl{wlc}] Tập trung cấu hình, quản lý tài nguyên vô tuyến (\ac{rrm}),
    chính sách bảo mật và chuyển vùng cho nhiều \ac{ap} dạng lightweight.
  \item[\acl{ds}] Đường trục (thường là Ethernet có dây) kết nối các \acp{ap}
    và cung cấp đường truy cập tới máy chủ cũng như Internet \cite{cisco_intro_80211, ciscopress_topologies}.
  \item[\acl{poe}] Các chuẩn 802.3af (~15~W), 802.3at (30~W) và 802.3bt
    (60--90~W) cấp nguồn cho \acp{ap} qua cáp Ethernet; \acp{ap} Wi-Fi~6/7
    hiện đại có thể cần 802.3at hoặc 802.3bt.
\end{description}

\subsection{Tập dịch vụ và cấu trúc mạng}

Chuẩn IEEE 802.11 định nghĩa các mô hình tập dịch vụ cơ bản nhằm phục vụ các ngữ cảnh triển khai khác nhau \cite{kurose2012, cisco_intro_80211, ciscopress_topologies}:
\begin{itemize}
  \item \textbf{\ac{bss}:} Một \ac{ap} cùng các trạm khách đã liên kết với nó. BSS
    được nhận diện duy nhất bằng \ac{bssid}, thông thường là địa chỉ MAC của cổng vô tuyến trên \ac{ap}.
  \item \textbf{\ac{ibss}:} Mạng ngang hàng (ad-hoc) trong đó các trạm giao tiếp trực tiếp với nhau mà không cần \ac{ap}.
  \item \textbf{\ac{ess}:} Nhiều \acp{bss} được kết nối liên thông qua hệ thống phân phối (\ac{ds}) và dùng chung một định danh \ac{ssid}, cho phép các thiết bị chuyển vùng liền mạch giữa các vùng phủ sóng.
  \item \textbf{Mesh \ac{bss}:} Các \acp{ap} chuyển tiếp lưu lượng không dây cho nhau qua nhiều chặng (theo chuẩn 802.11s hoặc giao thức riêng của nhà sản xuất); chỉ một số nút đóng vai trò cổng nối (mesh gateway) mới cần kết nối cáp vào mạng có dây.
\end{itemize}

\begin{figure}[htbp]
  \centering
  \begin{tikzpicture}[
      node distance=1.5cm and 1.0cm,
      ap/.style={draw,rounded corners,fill=blue!10,minimum width=2.4cm,minimum
      height=0.8cm,font=\small},
      sta/.style={draw,circle,fill=green!10,minimum
      size=0.9cm,font=\scriptsize},
      ds/.style={draw,fill=orange!15,minimum width=7.5cm,minimum
      height=0.7cm,font=\small}
    ]
    \node[ds] (ds) {Hệ thống phân phối (LAN có dây)};

    \node[ap, below left=1.2cm and 0.5cm of ds.center] (ap1) {AP$_1$};
    \node[ap, below right=1.2cm and 0.5cm of ds.center] (ap2) {AP$_2$};

    \node[sta, below left=1.0cm and 0.8cm of ap1] (s1) {STA};
    \node[sta, below right=1.0cm and 0.8cm of ap1] (s2) {STA};
    \node[sta, below left=1.0cm and 0.8cm of ap2] (s3) {STA};
    \node[sta, below right=1.0cm and 0.8cm of ap2] (s4) {STA};

    \draw (ap1.north) -- (ap1.north |- ds.south);
    \draw (ap2.north) -- (ap2.north |- ds.south);

    \foreach \s in {s1,s2}{\draw[dashed] (\s) -- (ap1);}
    \foreach \s in {s3,s4}{\draw[dashed] (\s) -- (ap2);}

    \draw[dashed,blue!60] ([yshift=-1.0cm]ap1.center) ellipse (2.6cm and 1.9cm);
    \draw[dashed,blue!60] ([yshift=-1.0cm]ap2.center) ellipse (2.6cm and 1.9cm);

    \node[font=\scriptsize,blue!60!black,below=1.7cm of ap1] {BSS$_1$};
    \node[font=\scriptsize,blue!60!black,below=1.7cm of ap2] {BSS$_2$};
  \end{tikzpicture}
  \caption{
    Tập dịch vụ mở rộng gồm hai BSS kết nối qua hệ thống phân phối có dây.
  }
  \label{fig:ess}
\end{figure}

\begin{table}[htbp]
  \centering\footnotesize
  \caption{So sánh giữa chế độ hạ tầng (Infrastructure Mode) và chế độ tùy biến (Ad-hoc Mode).}
  \label{tab:infra-vs-adhoc}
  \begin{tabularx}{\linewidth}{@{} L{0.65} L{1.15} L{1.2} @{}}
    \toprule
    \textbf{Tiêu chí} & \textbf{Chế độ hạ tầng (Infrastructure)} & \textbf{Chế độ tùy biến (Ad-hoc / IBSS)} \\
    \midrule
    Thiết bị trung tâm & Bắt buộc có \ac{ap} đóng vai trò điều phối trung tâm & Hoạt động ngang hàng, không có thiết bị trung tâm \\
    Cơ chế điều phối kênh & Có thể hỗ trợ DCF hoặc PCF/HCCA do \ac{ap} làm chủ & Phân tán hoàn toàn, chỉ hoạt động ở chế độ DCF (\ac{csmaca}) \\
    Cấp phát địa chỉ IP & Máy chủ DHCP tập trung trên mạng phân phối (\ac{ds}) & Gán tĩnh thủ công hoặc tự cấp IP nội bộ (APIPA / Link-local) \\
    Khả năng mở rộng & Rất tốt; dễ mở rộng thành \ac{ess} với nhiều \acp{bss} & Rất hạn chế; hiệu năng giảm mạnh khi số trạm tăng \\
    Bảo mật và kiểm soát & Toàn diện; hỗ trợ WPA2/WPA3, 802.1X/RADIUS doanh nghiệp & Hạn chế; chủ yếu dùng khóa chia sẻ trước (PSK) độc lập \\
    Ứng dụng thực tế & Mạng gia đình, văn phòng, trường học, doanh nghiệp lớn & Chia sẻ dữ liệu tạm thời, cứu nạn khẩn cấp, mạng cảm biến \\
    \bottomrule
  \end{tabularx}
\end{table}

\subsubsection{Các biến thể của kiến trúc tùy biến (Ad-hoc / IBSS)}

Kiến trúc \ac{ibss} trong chuẩn IEEE 802.11 cho phép các trạm \ac{sta} kết nối trực tiếp với nhau mà không cần sự hiện diện của một điểm truy cập trung tâm \cite{gast2005, kurose2012, cisco_intro_80211, hcmus_wlan_report}. Tùy thuộc vào phạm vi vô tuyến và cấu trúc kết nối, mạng Ad-hoc được phát triển thành các dạng biến thể chính:
\begin{itemize}
  \item \textbf{Ad-hoc một chặng (Single-hop Ad-hoc):} Tất cả các nút mạng bắt buộc phải nằm trong tầm phủ sóng vô tuyến trực tiếp của nhau. Mọi khung tin được gửi trực tiếp từ trạm nguồn đến trạm đích. Dạng mạng này có cấu hình đơn giản nhưng phạm vi hoạt động bị giới hạn trong bán kính phủ sóng của một ô duy nhất.
  \item \textbf{Ad-hoc nhiều chặng (\acl{lan} di động tùy biến -- MANET):} Khi các trạm nằm ngoài tầm phủ sóng trực tiếp của nhau, các nút trung gian đóng vai trò như các bộ định tuyến để chuyển tiếp gói tin qua nhiều chặng vô tuyến. Việc định tuyến trong MANET đòi hỏi các giao thức chuyên dụng, chia thành hai nhóm chính:
    \begin{itemize}[leftmargin=1.5em]
      \item \emph{Giao thức định tuyến phản ứng (Reactive / On-demand):} Điển hình là AODV (Ad hoc On-Demand Distance Vector) và DSR (Dynamic Source Routing), chỉ tiến hành tìm kiếm và thiết lập tuyến đường khi trạm nguồn có nhu cầu gửi dữ liệu, giúp tiết kiệm băng thông khi mạng ít truyền tin.
      \item \emph{Giao thức định tuyến chủ động (Proactive / Table-driven):} Điển hình là OLSR (Optimized Link State Routing), định kỳ trao đổi thông tin trạng thái liên kết để liên tục duy trì bảng định tuyến toàn cục tại mỗi nút, giúp truyền dữ liệu ngay lập tức mà không mất trễ dò đường.
    \end{itemize}
  \item \textbf{Mô hình kết hợp (Hybrid) và Wi-Fi Direct:} Để khắc phục nhược điểm phức tạp và thiếu tương thích của IBSS truyền thống, chuẩn Wi-Fi Direct (Wi-Fi P2P) đưa ra cơ chế kết nối ngang hàng cải tiến: các thiết bị tự động thương lượng để một trạm đảm nhận vai trò Trưởng nhóm (P2P Group Owner -- P2P GO), hoạt động như một \ac{ap} ảo cung cấp dịch vụ DHCP và quản lý liên kết, trong khi các trạm khác kết nối với tư cách P2P Client.
\end{itemize}

\subsubsection{Hệ thống phân phối không dây (WDS) và cơ chế bắc cầu 4 địa chỉ MAC}

Thông thường, hệ thống phân phối (\ac{ds}) được triển khai bằng cáp mạng Ethernet có dây. Tuy nhiên, trong các địa hình khó đi dây hoặc cần mở rộng vùng phủ sóng nhanh, cơ chế \ac{ds} không dây (Wireless Distribution System -- WDS) cho phép các \acp{ap} liên kết trực tiếp với nhau qua sóng vô tuyến để tạo thành một mạng xương sống không dây \cite{gast2005, coleman2021}.

Trong mạng 802.11 thông thường giữa trạm khách và \ac{ap}, khung dữ liệu chỉ sử dụng 3 trường địa chỉ MAC (trong đó các cờ To DS và From DS không đồng thời bằng 1). Ngược lại, WDS yêu cầu sử dụng đầy đủ cả 4 trường địa chỉ MAC trong tiêu đề khung 802.11 bằng cách thiết lập đồng thời cả hai cờ To DS = 1 và From DS = 1 trong trường Frame Control:
\begin{itemize}
  \item \textbf{Địa chỉ 1 (Receiver Address -- RA):} Địa chỉ MAC của \ac{ap} nhận khung vô tuyến ở chặng hiện tại.
  \item \textbf{Địa chỉ 2 (Transmitter Address -- TA):} Địa chỉ MAC của \ac{ap} phát khung vô tuyến ở chặng hiện tại.
  \item \textbf{Địa chỉ 3 (Destination Address -- DA):} Địa chỉ MAC của trạm đích nhận dữ liệu cuối cùng trong mạng.
  \item \textbf{Địa chỉ 4 (Source Address -- SA):} Địa chỉ MAC của trạm nguồn phát sinh gói dữ liệu ban đầu.
\end{itemize}
Cơ chế 4 địa chỉ này cho phép bắc cầu nguyên vẹn khung mạng Ethernet lớp 2 qua môi trường vô tuyến mà không làm thay đổi thông tin địa chỉ gốc của thiết bị đầu cuối, làm nền tảng cho các bộ lặp không dây (wireless repeater) và kiến trúc mạng lưới Mesh 802.11s.

\subsection{Mô hình triển khai}

\begin{table}[htbp]
  \centering\footnotesize
  \caption{Các mô hình triển khai WLAN phổ biến.}
  \label{tab:deployment-models}
  \begin{tabularx}{\linewidth}{@{} L{0.8} L{0.8} L{0.8} L{1.3} L{1.3} @{}}
    \toprule
    \textbf{Mô hình} & \textbf{Mặt phẳng điều khiển} & \textbf{Mặt phẳng dữ liệu} & \textbf{Ưu điểm} & \textbf{Hạn chế} \\
    \midrule
    \acp{ap} tự chủ (độc lập) & Từng \ac{ap} & Chuyển mạch cục bộ & Đơn giản, tiết kiệm khi chỉ có 1--3 \acp{ap} & Khó mở rộng và quản lý đồng nhất \\
    Dùng bộ điều khiển (tập trung) & \acs{wlc} & Tạo đường hầm tới \acs{wlc} qua \acs{capwap} & Chính sách tập trung, chuyển vùng lớp 3 thuận tiện & Bộ điều khiển có thể thành nút nghẽn \\
    Dùng bộ điều khiển (chuyển mạch cục bộ) & \acs{wlc} & Chuyển mạch tại \ac{ap} (ví dụ FlexConnect) & Phù hợp chi nhánh, vẫn hoạt động khi mất WAN & Thiết kế VLAN phức tạp hơn \\
    Quản lý qua đám mây & Dịch vụ đám mây & Cục bộ tại \ac{ap} & Cấp cấu hình tự động, có phân tích & Quản lý phụ thuộc Internet \\
    Mesh & Phân tán hoặc qua bộ điều khiển & Không dây nhiều chặng & Không cần kéo cáp tới mọi \ac{ap} & Thông lượng giảm qua mỗi chặng \\
    \bottomrule
  \end{tabularx}
\end{table}

\subsubsection{Kiến trúc Split-MAC và giao thức CAPWAP}

Trong mô hình tập trung sử dụng bộ điều khiển, kiến trúc phân chia chức năng MAC (Split-MAC Architecture) được áp dụng nhằm tối ưu hóa hiệu năng và đơn giản hóa công tác quản trị \cite{coleman2021, rfc5416}. Chức năng của lớp MAC được phân định rõ rệt giữa hai thành phần:
\begin{itemize}
  \item \textbf{Real-time MAC (tác vụ thời gian thực tại Lightweight AP):} Đảm nhiệm các chức năng vô tuyến nhạy cảm với thời gian đòi hỏi phản hồi tức thì ở mức mili-giây, bao gồm: phát khung báo hiệu (Beacon), gửi phản hồi dò quét (Probe Response), truyền khung xác nhận (ACK), cơ chế bắt tay bảo lưu kênh CSMA/CA, phát hiện sóng mang, đệm khung dữ liệu cục bộ và thực hiện mã hóa/giải mã phần cứng vô tuyến (như CCMP/GCMP).
  \item \textbf{Non-real-time MAC (tác vụ phi thời gian thực tại WLC):} Đảm nhiệm các chức năng quản lý logic và xử lý tập trung, bao gồm: xác thực người dùng 802.1X/EAP và chuyển tiếp RADIUS, quản lý tài nguyên vô tuyến (\ac{rrm} -- tự động điều phối công suất phát và phân bổ kênh tránh nhiễu), điều phối chuyển vùng (roaming) mượt mà giữa các AP, thực thi chính sách kiểm soát truy cập/QoS và quản trị cấu hình cũng như firmware tập trung.
\end{itemize}

Sự phân chia này giúp các \acp{ap} trở nên nhẹ tải, chi phí phần cứng thấp hơn, đồng thời toàn bộ chính sách và thông tin điều khiển mạng được quản lý đồng nhất tại \ac{wlc}. Giao thức \ac{capwap} \cite{rfc5415, rfc5416} cung cấp đường hầm liên lạc giữa AP và WLC, bao gồm kênh điều khiển (cổng UDP 5246, bắt buộc mã hóa bảo vệ bằng DTLS) và kênh dữ liệu (cổng UDP 5247, tùy chọn mã hóa).

\subsection{Các dịch vụ mạng hỗ trợ WLAN}

Trong mô hình OSI/TCP-IP, WLAN chủ yếu hoạt động ở lớp vật lý (Physical Layer,
PHY) và lớp liên kết dữ liệu (Data Link Layer), theo các đặc tả PHY/MAC của
IEEE 802.11. Lớp MAC điều phối quyền truy cập môi trường vô tuyến dùng chung;
các lớp mạng và lớp cao hơn tiếp tục đảm nhiệm định tuyến IP, vận chuyển và
các ứng dụng như trên mạng có dây. Vì vậy, một thiết bị WLAN vẫn cần các dịch
vụ nền như DHCP để nhận cấu hình IP và DNS để phân giải tên miền.

\begin{itemize}
  \item \textbf{VLAN và SSID:} mỗi \ac{ssid} thường được ánh xạ tới một
    \ac{vlan} (doanh nghiệp, khách, thoại, \ac{iot}) để phân đoạn và kiểm soát
    lưu lượng.
  \item \textbf{DHCP và DNS:} cần được cấu hình đủ năng lực cho số lượng thiết
    bị cực đại; mạng khách thường dùng thời hạn thuê địa chỉ ngắn.
  \item \textbf{RADIUS:} cung cấp xác thực, phân quyền và ghi nhận hoạt động
    cho 802.1X (Mục~\ref{sec:security}).
  \item \textbf{Chuyển vùng lớp 2 và lớp 3:} trong cùng một mạng con, thiết bị
    giữ nguyên địa chỉ IP; khi chuyển giữa các mạng con, bộ điều khiển có thể
    tạo đường hầm lưu lượng tới điểm neo hoặc thiết bị phải nhận địa chỉ IP mới
    qua DHCP.
\end{itemize}

Các báo cáo tham khảo cũng minh họa cách WLAN được tích hợp với VLAN, DHCP,
DNS và các dịch vụ mạng doanh nghiệp; nội dung cấu hình và kết quả kiểm thử
không được đưa vào phần lý thuyết này \cite{hcmus_office_network}.
Trên phương diện thiết bị đầu cuối, các máy trạm máy tính chạy Windows \cite{microsoft_wifi} hay thiết bị di động chạy Android \cite{google_android_wifi} và iOS \cite{apple_ios_wifi} đều tích hợp trình quản lý mạng không dây để dò quét SSID, xác thực bảo mật và tiếp nhận cấu hình IP từ DHCP một cách liền mạch.

\begin{tomtat}
\item Thành phần cơ bản: \acs{sta} (trạm), \acs{ap} (điểm truy cập), \acs{wlc} (bộ điều khiển),
  \acs{ds} (hệ thống phân phối) và \acs{poe} (cấp nguồn qua cáp Ethernet).
\item Các tập dịch vụ: \acs{bss} (hạ tầng với một \acs{ap}), \acs{ibss} (ad-hoc ngang hàng),
  \acs{ess} (nhiều \acs{bss} kết nối qua \acs{ds}) và Mesh \acs{bss}.
\item Kiến trúc Ad-hoc mở rộng gồm: Single-hop (một chặng), Multi-hop MANET (với các giao thức
  định tuyến phản ứng AODV/DSR và chủ động OLSR), cùng mô hình lai Wi-Fi Direct (P2P Group Owner).
\item Cơ chế WDS (Wireless Distribution System) dùng tiêu đề khung 802.11 đủ 4 địa chỉ MAC (To DS=1,
  From DS=1) để bắc cầu các AP không dây mà vẫn giữ nguyên địa chỉ nguồn/đích lớp 2.
\item Mô hình triển khai: \acs{ap} độc lập, dùng bộ điều khiển (tập trung hoặc chuyển mạch cục bộ),
  quản lý qua đám mây và mesh.
\item Cơ chế Split-MAC phân chia: Real-time MAC (beacon, ACK, CSMA/CA, mã hóa vô tuyến) xử lý tại AP;
  Non-real-time MAC (802.1X/RADIUS, RRM, chuyển vùng, chính sách) quản trị tập trung tại \acs{wlc}.
\item Giao thức \acs{capwap} (UDP 5246 cho điều khiển qua DTLS, 5247 cho dữ liệu)
  kết nối và quản lý các \acs{ap} lightweight từ \acs{wlc}.
\item Dịch vụ xung quanh: ánh xạ \acs{ssid}--\acs{vlan} để phân
  đoạn lưu lượng, DHCP/DNS, \acs{radius} cho 802.1X, và chuyển vùng lớp 2/lớp 3.
\end{tomtat}
\endgroup
